\begin{lemma}[The two budgets of a stopped contact walk]
Let $F$ be a finite-dimensional real inner-product space, let
$1\le C_0<\infty$, let $\kappa>0$ and $H\in\R$, and let
$v\in\Aclass_{m,C_0}$ in the sense of \noderef{IsFitClass} be centered at
$z_*=(0,1)$ with constant $C_0$ in the sense of \noderef{IsCenteredAt} and
satisfy the horizontal Jensen bound with constant $\kappa$ at every time in the
sense of \noderef{HasJensenBound}. Let $W$ be the contact envelope of
\noderef{ContactEnv} formed from $v$ and $\kappa$, and assume that $W$ is
continuous on $\{(x,t):t>0\}$ and has subquadratic spatial growth on compact
subintervals of $(0,\infty)$ in the sense of \noderef{SubquadraticGrowth}.
Write $J_\epsilon$ for the quantity of \noderef{JGap}, $\delta_j=2^{-j}$ and
$L(K)=1+\log K$.

Then there is a constant $C_*>0$ such that for every integer $T\ge2$, every
real $R\ge2$, every integer $N\ge1$, every $\epsilon$, every measure space
$(\Omega,\mathcal M,\mu)$ with a filtration $(\mathcal G_n)$, and all functions
$\xi_n$, $z_n$, $z^+_n$, $z^-_n$, $Q_0(n)$ and $\tau$ on it which form a stopped
contact walk for $(\kappa,H,T,R,v,N,\epsilon)$ in the sense of
\noderef{IsContactWalk}, the following hold, where
\[
  Q_+(n)=\frac{x^+_n-x_n}{\epsilon},\qquad
  Q_-(n)=\frac{x_n-x^-_n}{\epsilon},\qquad
  \mathcal E_n=|Q_+(n)-Q_0(n)|^2+|Q_-(n)-Q_0(n)|^2,
\]
and $B=R+T+2$:
\begin{enumerate}
\item[(i)] for every stopping time $\sigma$ for $(\mathcal G_n)$ with
$\sigma\le\tau$ pointwise, the function $J_\epsilon(z_\sigma)$ is
$\mu$-integrable and $\mathbb E J_\epsilon(z_\sigma)\ge J_\epsilon(z_*)$;
\item[(ii)] the function $\sum_{n<\tau}\mathcal E_n$ is $\mu$-integrable and
$\mathbb E\sum_{n<\tau}\mathcal E_n\le C_*L(B)^4$.
\end{enumerate}
\end{lemma}
\begin{proof}
\emph{Step 0: the constants.}
Let $C_1>0$ be the constant of \noderef{FittedEnvelopeBounds} for this $C_0$
and put
\[
  c=\frac{1}{16(1+\log2)}>0,\qquad C_*=\frac{C_1+|H|+1}{c}>0 ,
\]
so that $c$ is the constant appearing in \noderef{DPDrift}. Both depend only on
$C_0$, $\dim F$ and $H$, hence neither depends on $T$, $R$, $N$ or on the walk.
Fix now $T$, $R$, $N$, $\epsilon$ and a stopped contact walk as in the
statement, and write $\delta=\delta_N$, $B=R+T+2\ge6$ and $L=L(B)\ge1$. By
clauses 2 and 3 of \noderef{IsContactWalk} we have $0<\epsilon\le\delta/4$, so
$\epsilon\le1/8$ and $\epsilon\le B$, and $\mu$ is a probability measure by
clause 1.

\emph{Step 1: the walk stays in a compact box.}
We claim that for every $n$ and every $\omega$
\begin{equation}\label{eq:cwb-box}
  \delta-\epsilon\le t_n\le T+\epsilon,\qquad |x_n|\le R+1 ,
\end{equation}
where $z_n=(x_n,t_n)$. For the time bounds we induct on $n$. At $n=0$, clause 6
gives $z_0=z_*$, so $t_0=1$, and $\delta\le1/2<1\le T$ because $N\ge1$ and
$T\ge2$. Assume the bounds at $n$. If $n\ge\tau(\omega)$ then clause 7 gives
$z_{n+1}=z_n$ and there is nothing to prove. If $n<\tau(\omega)$ then clause 7
gives $z_{n+1}\in\{z^+_n,z^-_n\}$, and by clause 6 together with
\noderef{IsStepPair} the time coordinates of $z^\pm_n$ are $t_n\pm\epsilon$;
clause 10 gives $\delta<t_n<T$, so $t_{n+1}\ge t_n-\epsilon>\delta-\epsilon$ and
$t_{n+1}\le t_n+\epsilon<T+\epsilon$. In particular $t_n\ge\delta-\epsilon>0$
for every $n$, since $\epsilon\le\delta/4$. For the spatial bound: clause 12
gives $|x_n|\le R+1$ for $n\le\tau$, and for $n>\tau$ clause 7 gives
$z_n=z_\tau$ by induction, so the same bound holds.

\emph{Step 2: two-sided bounds for $J_\epsilon$ on the box.}
Let $w=(y,s)$ satisfy $|y|\le R+1$ and $\delta-\epsilon\le s\le T+\epsilon$;
then $s>0$ and $|s|\le T+1\le B$ and $|y|\le R+1\le B$.

\emph{The upper bound.} By \noderef{JUpperBound}, applied with $\epsilon>0$ and
$s>0$, $J_\epsilon(w)\le\tfrac12(\Senv^+_\epsilon v(w)+\Senv^-_\epsilon v(w))$
with the envelope slopes of \noderef{EnvSlopeP} and \noderef{EnvSlopeM}; and by
the first clause of \noderef{FittedEnvelopeBounds}, applied to $v$ with the
centering point $(0,1)$, the box size $K=B\ge2$ and the length
$\epsilon\le B$, both slopes are at most $C_1L^4$. Hence
\begin{equation}\label{eq:cwb-Jup}
  J_\epsilon(w)\le C_1L^4 .
\end{equation}

\emph{The lower bound.} The standing hypotheses on $v$, on $W$ and on the
scales are exactly those of \noderef{JGapLowerBoundAndContinuity}, applied with
the spatial radius $R+1$, the time $T$, the scale $\delta$ and the length
$\epsilon$: here $0\le R+1$ and $0\le T$ because $R\ge2$ and $T\ge2$, while
$0<\delta\le1/2$ and $0<\epsilon\le\delta/4$ hold by Step 0. Its first
conclusion supplies a constant $\Lambda_0\ge0$ with
\begin{equation}\label{eq:cwb-Jlow}
  J_\epsilon(w)\ge-\Lambda_0
\end{equation}
for every $w=(y,s)$ with $|y|\le R+1$ and $\delta-\epsilon\le s\le T+\epsilon$,
which is every $w$ as above. Put
$\Lambda_1=\max\{C_1L^4,\Lambda_0\}$; by \eqref{eq:cwb-Jup},
\eqref{eq:cwb-Jlow} and Step 1, $|J_\epsilon(z_n)|\le\Lambda_1$ for every $n$
and $\omega$.

\emph{Step 3: measurability.}
Let $\rho(y,s)=(y,\max\{s,\delta/2\})$. By the second conclusion of
\noderef{JGapLowerBoundAndContinuity}, applied with the same data, the
composite $J_\epsilon\circ\rho$ is continuous on all of space-time, and it
agrees with $J_\epsilon$ at every point whose time coordinate is at least
$\delta/2$, because there $\max\{s,\delta/2\}=s$ --- in particular, by Step 1
and
$\epsilon\le\delta/4$, at every $z_n(\omega)$. Since each $z_n$ is
$\mathcal G_n$-measurable by clause 4 of \noderef{IsContactWalk}, the function
$\omega\mapsto J_\epsilon(z_n(\omega))=(J_\epsilon\circ\rho)(z_n(\omega))$ is
$\mathcal G_n$-measurable, being a continuous function of a
$\mathcal G_n$-measurable one; the same argument applies to $z^\pm_n$ with
$J_\epsilon\circ\rho$ in place of $J_\epsilon$. Likewise each $\mathcal E_n$ is
$\mathcal G_n$-measurable, being a continuous function of the
$\mathcal G_n$-measurable $z_n$, $z^\pm_n$ and $Q_0(n)$, and
$\mathcal E_n\ge0$, being a sum of two squares. Finally, $\tau$ is a stopping
time by clause 9, so $\{n<\tau\}\in\mathcal G_n$ for every $n$, and $\tau$ is
measurable because $\{\tau=i\}\in\mathcal G_i$ for every $i$ and $\tau$ has
countable range.

\emph{Step 4: almost everywhere the step states lie in the same box.}
Fix $n$ and let $G$ be the set of $\omega$ with
$|x^+_n(\omega)|\le R+1$ and $\delta-\epsilon\le t^+_n(\omega)\le T+\epsilon$,
a $\mathcal G_n$-measurable set by Step 3. Put
$S=\{n<\tau\}\setminus G\in\mathcal G_n$. On $S\cap\{\xi_n=1\}$ clause 7 gives
$z_{n+1}=z^+_n$, so \eqref{eq:cwb-box} applied at $n+1$ puts $\omega$ in $G$;
hence $S\cap\{\xi_n=1\}=\emptyset$. On the other hand $\xi_n$ is bounded and
measurable, hence integrable, and clause 3 of \noderef{IsContactWalk} gives
$\mathbb E[\xi_n\mid\mathcal G_n]=0$ a.e., so by the characterizing property of
the conditional expectation $\int_S\xi_n\,d\mu=0$. But $\xi_n\in\{1,-1\}$ by
clause 3, so $\xi_n=-1$ on $S$ and therefore
$\int_S\xi_n\,d\mu=-\mu(S)$. Hence $\mu(S)=0$, that is, almost everywhere on
$\{n<\tau\}$ the state $z^+_n$ satisfies the bounds of \eqref{eq:cwb-box}. The
same argument with the sign $-1$ in place of $+1$ --- on
$S'\cap\{\xi_n=-1\}$ clause 7 gives $z_{n+1}=z^-_n$, and on $S'$ we then have
$\xi_n=+1$, so $\int_{S'}\xi_n\,d\mu=\mu(S')$ --- gives the same conclusion for
$z^-_n$. Consequently, by Step 2, almost everywhere on $\{n<\tau\}$
\begin{equation}\label{eq:cwb-stepbd}
  |J_\epsilon(z^+_n)|\le\Lambda_1,\qquad |J_\epsilon(z^-_n)|\le\Lambda_1 .
\end{equation}

\emph{Step 5: the one-step drift.}
Let $n<\tau(\omega)$. Then clause 10 gives $t_n>\delta\ge4\epsilon>\epsilon$,
clause 6 says that $(z^+_n,z^-_n)$ is an optimizing step pair of size $\epsilon$
at $z_n$ in the sense of \noderef{IsStepPair}, and clause 8 says that
$\Penv_\epsilon W(\cdot,t_n)$ has gradient $Q_0(n)$ at $x_n$. So
\noderef{DPDrift} applies at $z_n$ with this step pair and this reference
velocity, with the four gaps $D,e,\tilde D,\tilde e$ given by its defining
equations. Its sixth and seventh conclusions read
\[
  \frac{D+e}{4\epsilon}
   \le\frac{J_\epsilon(z^+_n)+J_\epsilon(z^-_n)}{2}-J_\epsilon(z_n),
  \qquad
  c\,\mathcal E_n\le\frac{D+e}{4\epsilon},
\]
and together
\begin{equation}\label{eq:cwb-drift}
  c\,\mathcal E_n
   \le\frac{J_\epsilon(z^+_n)+J_\epsilon(z^-_n)}{2}-J_\epsilon(z_n)
  \qquad(n<\tau).
\end{equation}
With \eqref{eq:cwb-stepbd} and $|J_\epsilon(z_n)|\le\Lambda_1$ this gives the
a.e.\ bound
\begin{equation}\label{eq:cwb-Ebd}
  c\,\mathcal E_n\le\Lambda:=C_1L^4+\Lambda_0
  \qquad\text{a.e.\ on }\{n<\tau\} ,
\end{equation}
and since the index set is countable, almost every $\omega$ satisfies
\eqref{eq:cwb-Ebd} for every $n<\tau(\omega)$ simultaneously.

\emph{Step 6: the drift submartingale and optional stopping.}
For $n\ge0$ put
\[
  M_n=J_\epsilon(z_n)-c\sum_{k<n}\mathbf1_{\{k<\tau\}}\mathcal E_k .
\]
Each $M_n$ is $\mathcal G_n$-measurable by Step 3, and integrable: $\mu$ is a
probability measure, $|J_\epsilon(z_n)|\le\Lambda_1$ by Step 2, and by
\eqref{eq:cwb-Ebd} the sum has at most $n$ terms each at most $\Lambda/c$
almost everywhere.

We claim $(M_n)$ is a submartingale. Set
\[
  A_n=\mathbf1_{\{n<\tau\}}
      \Bigl(\frac{J_\epsilon(z^+_n)+J_\epsilon(z^-_n)}{2}-J_\epsilon(z_n)
            -c\,\mathcal E_n\Bigr),\qquad
  D_n=\mathbf1_{\{n<\tau\}}\,
      \frac{J_\epsilon(z^+_n)-J_\epsilon(z^-_n)}{2},
\]
both $\mathcal G_n$-measurable by Step 3 (and both read through
$J_\epsilon\circ\rho$, which is where their measurability comes from). Then
\begin{equation}\label{eq:cwb-decomp}
  M_{n+1}-M_n=A_n+\xi_nD_n .
\end{equation}
Indeed, on $\{n\ge\tau\}$ clause 7 gives $z_{n+1}=z_n$ and both sides vanish; on
$\{n<\tau\}$ clause 7 gives $z_{n+1}=z^+_n$ where $\xi_n=1$ and
$z_{n+1}=z^-_n$ where $\xi_n=-1$, and since $\xi_n\in\{1,-1\}$ by clause 3,
\[
  J_\epsilon(z_{n+1})
   =\frac{J_\epsilon(z^+_n)+J_\epsilon(z^-_n)}{2}
    +\xi_n\frac{J_\epsilon(z^+_n)-J_\epsilon(z^-_n)}{2}
\]
in both cases, which is \eqref{eq:cwb-decomp} after subtracting
$J_\epsilon(z_n)+c\,\mathcal E_n$. By \eqref{eq:cwb-drift}, $A_n\ge0$ almost
everywhere. By \eqref{eq:cwb-stepbd}, $D_n$ is bounded by $\Lambda_1$ almost
everywhere, so the pull-out property of the conditional expectation for a
bounded $\mathcal G_n$-measurable factor and clause 3 give
$\mathbb E[\xi_nD_n\mid\mathcal G_n]=D_n\,\mathbb E[\xi_n\mid\mathcal G_n]=0$
a.e. Since $A_n$ is $\mathcal G_n$-measurable and integrable,
$\mathbb E[A_n\mid\mathcal G_n]=A_n$ a.e., and therefore
$\mathbb E[M_{n+1}-M_n\mid\mathcal G_n]=A_n\ge0$ almost everywhere. As each
$M_n$ is adapted and integrable, $(M_n)$ is a submartingale.

Let $\sigma$ be a stopping time with $\sigma\le\tau$ and let $h\in\mathbb N$.
The constant $0$ and $\sigma\wedge h$ are stopping times with
$0\le\sigma\wedge h\le h$, so optional stopping for submartingales gives
$\mathbb E M_0\le\mathbb E M_{\sigma\wedge h}$. Now $M_0=J_\epsilon(z_*)$ by
clause 6, and for $k<\sigma\wedge h\le\tau$ the indicators in the sum equal $1$,
so
\begin{equation}\label{eq:cwb-trunc}
  J_\epsilon(z_*)+c\,\mathbb E\sum_{k<\sigma\wedge h}\mathcal E_k
   \le\mathbb E J_\epsilon(z_{\sigma\wedge h}) ,
\end{equation}
both expectations existing: the integrand on the right is bounded by
$\Lambda_1$, and the sum on the left is nonnegative and bounded by
$h\Lambda/c$ almost everywhere, while both are measurable because
$\sigma\wedge h\le h$ makes them finite sums of the measurable
$\mathbf1_{\{\sigma\wedge h=i\}}$ times measurable functions.

\emph{Step 7: $h\to\infty$, and the two conclusions.}
Fix $\sigma$ as above. For every $\omega$ we have $\sigma\wedge h=\sigma$ as
soon as $h\ge\sigma(\omega)$, so $J_\epsilon(z_{\sigma\wedge h})\to
J_\epsilon(z_\sigma)$ pointwise and $\sum_{k<\sigma\wedge h}\mathcal E_k$
increases to $\sum_{k<\sigma}\mathcal E_k$ pointwise. In particular
$J_\epsilon(z_\sigma)$ is measurable, being a pointwise limit of measurable
functions, and $|J_\epsilon(z_\sigma)|\le\Lambda_1$ by Step 2, so it is
integrable: this is the first half of (i). Dominated convergence, legitimate
because all these functions are bounded by the constant $\Lambda_1$ and $\mu$ is
a probability measure, gives
$\mathbb E J_\epsilon(z_{\sigma\wedge h})\to\mathbb E J_\epsilon(z_\sigma)$.
Dropping the nonnegative error term in \eqref{eq:cwb-trunc} leaves
$J_\epsilon(z_*)\le\mathbb E J_\epsilon(z_{\sigma\wedge h})$ for every $h$, and
letting $h\to\infty$ gives $J_\epsilon(z_*)\le\mathbb E J_\epsilon(z_\sigma)$.
This is (i).

For (ii) take $\sigma=\tau$ in \eqref{eq:cwb-trunc}. By \eqref{eq:cwb-Jup} and
Step 1, $\mathbb E J_\epsilon(z_{\tau\wedge h})\le C_1L^4$, and by clause 13 of
\noderef{IsContactWalk}, $J_\epsilon(z_*)\ge H-1$. Hence for every $h$
\[
  c\,\mathbb E\sum_{k<\tau\wedge h}\mathcal E_k
   \le C_1L^4-(H-1)\le\bigl(C_1+|H|+1\bigr)L^4 ,
\]
using $L\ge1$ and $-H\le|H|$, that is
$\mathbb E\sum_{k<\tau\wedge h}\mathcal E_k\le C_*L^4$. The functions
$\sum_{k<\tau\wedge h}\mathcal E_k$ are nonnegative, measurable and
nondecreasing in $h$ with pointwise limit $\sum_{k<\tau}\mathcal E_k$, so by
monotone convergence --- applied to the integrals with values in $[0,\infty]$,
which is what makes the conclusion carry integrability with it ---
\[
  \int_\Omega\sum_{k<\tau}\mathcal E_k\,d\mu
   =\sup_h\int_\Omega\sum_{k<\tau\wedge h}\mathcal E_k\,d\mu\le C_*L^4<\infty .
\]
The function $\sum_{k<\tau}\mathcal E_k$ is therefore integrable, with
$\mathbb E\sum_{k<\tau}\mathcal E_k\le C_*L^4$. This is (ii). \end{proof}
